\documentclass[11pt]{article}
\usepackage[utf8]{inputenc}
\usepackage[T1]{fontenc}
\usepackage{lmodern}
% jugeo-common.tex — shared preamble for all Judgment Geometry papers
% Include via: % jugeo-common.tex — shared preamble for all Judgment Geometry papers
% Include via: % jugeo-common.tex — shared preamble for all Judgment Geometry papers
% Include via: \input{jugeo-common}

% ── Packages ────────────────────────────────────────────────────────────
\usepackage{amsmath,amssymb,amsthm,mathtools}
\usepackage{tikz-cd}
\usepackage{listings}
\usepackage{xcolor}
\usepackage{xspace}
\usepackage{booktabs}
\usepackage{multirow}
\usepackage{enumitem}
\usepackage[expansion=false]{microtype}
\usepackage{hyperref}
\usepackage{cleveref}
% \usepackage{thmtools}  % disabled: not available in this TeX installation
\usepackage{float}
\usepackage{caption}
\usepackage{subcaption}
\usepackage{tabularx}
\usepackage{stmaryrd}       % \llbracket, \rrbracket
\usepackage{bbm}            % \mathbbm{1}

% ── Page geometry ───────────────────────────────────────────────────────
\usepackage[margin=1in]{geometry}

% ── Theorem environments ────────────────────────────────────────────────
\theoremstyle{plain}
\newtheorem{theorem}{Theorem}[section]
\newtheorem{lemma}[theorem]{Lemma}
\newtheorem{proposition}[theorem]{Proposition}
\newtheorem{corollary}[theorem]{Corollary}

\theoremstyle{definition}
\newtheorem{definition}[theorem]{Definition}
\newtheorem{example}[theorem]{Example}
\newtheorem{construction}[theorem]{Construction}

\theoremstyle{remark}
\newtheorem{remark}[theorem]{Remark}
\newtheorem{notation}[theorem]{Notation}

% ── Python listings style ──────────────────────────────────────────────
\definecolor{jugeoBlue}{HTML}{2563EB}
\definecolor{jugeoGreen}{HTML}{059669}
\definecolor{jugeoRed}{HTML}{DC2626}
\definecolor{jugeoPurple}{HTML}{7C3AED}
\definecolor{jugeoGray}{HTML}{6B7280}
\definecolor{jugeoBg}{HTML}{F8FAFC}

\lstdefinestyle{jugeo-python}{
  language=Python,
  basicstyle=\ttfamily\small,
  keywordstyle=\color{jugeoBlue}\bfseries,
  stringstyle=\color{jugeoGreen},
  commentstyle=\color{jugeoGray}\itshape,
  numberstyle=\tiny\color{jugeoGray},
  backgroundcolor=\color{jugeoBg},
  numbers=left,
  numbersep=6pt,
  frame=single,
  framerule=0.4pt,
  rulecolor=\color{jugeoGray!40},
  breaklines=true,
  breakatwhitespace=true,
  showstringspaces=false,
  tabsize=4,
  xleftmargin=1.5em,
  framexleftmargin=1.5em,
  morekeywords={dataclass,frozen,field,Enum,IntEnum,auto,Any,
                Mapping,Sequence,Optional,match,case},
  literate={→}{{$\to$}}1 {←}{{$\leftarrow$}}1
           {≤}{{$\leq$}}1 {≥}{{$\geq$}}1 {≠}{{$\neq$}}1
           {∈}{{$\in$}}1 {∀}{{$\forall$}}1 {∃}{{$\exists$}}1
           {⊕}{{$\oplus$}}1 {⊖}{{$\ominus$}}1,
  escapeinside={(*@}{@*)},
}
\lstset{style=jugeo-python}

\lstdefinestyle{jugeo-output}{
  basicstyle=\ttfamily\small,
  backgroundcolor=\color{jugeoBg},
  frame=single,
  framerule=0.4pt,
  rulecolor=\color{jugeoGray!40},
  breaklines=true,
  showstringspaces=false,
  xleftmargin=1.5em,
  framexleftmargin=0.5em,
  numbers=none,
}

% ── JuGeo-specific macros ──────────────────────────────────────────────

% Judgment 8-tuple
\newcommand{\J}{\mathcal{J}}
\newcommand{\Jud}[8]{(#1,\,#2,\,#3,\,#4,\,#5,\,#6,\,#7,\,#8)}
\newcommand{\JudShort}[1]{\J_{#1}}

% Components
\newcommand{\coord}{\mathsf{c}}            % coordinate
\newcommand{\prop}{\varphi}                 % proposition
\newcommand{\carrier}{\mathcal{A}}          % carrier type
\newcommand{\evid}{\mathcal{E}}             % evidence bundle
\newcommand{\oblig}{\mathcal{O}}            % obligations
\newcommand{\obstr}{\mathcal{B}}            % obstructions
\newcommand{\trust}{\mathcal{T}}            % trust annotation
\newcommand{\prov}{\Pi}                     % provenance

% Trust algebra
\newcommand{\Talg}{\mathfrak{T}}            % trust ordered algebra
\newcommand{\Eadm}{\mathcal{E}_{\mathrm{adm}}}
\newcommand{\tleq}{\preceq}                 % trust ordering
\newcommand{\tjoin}{\oplus}                 % conservative join
\newcommand{\tatten}{\ominus}               % attenuation
\newcommand{\tpromote}{\uparrow_\pi}        % promotion
\newcommand{\tdemote}{\downarrow_\chi}       % demotion

% Trust tiers
\newcommand{\Tcontra}{\ensuremath{\mathsf{CONTRADICTED}}}
\newcommand{\Tunver}{\ensuremath{\mathsf{UNVERIFIED}}}
\newcommand{\Tcopilot}{\ensuremath{\mathsf{COPILOT}}}
\newcommand{\Toracle}{\ensuremath{\mathsf{ORACLE}}}
\newcommand{\Truntime}{\ensuremath{\mathsf{RUNTIME}}}
\newcommand{\Tsolver}{\ensuremath{\mathsf{SOLVER}}}
\newcommand{\Tproof}{\ensuremath{\mathsf{PROOF}}}

% Site and geometry
\newcommand{\Site}{\mathcal{S}}             % semantic site
\newcommand{\Cov}{\mathrm{Cov}}            % covering family
\newcommand{\Ob}{\mathrm{Ob}}              % objects
\newcommand{\Mor}{\mathrm{Mor}}            % morphisms
\newcommand{\res}{\mathrm{res}}            % restriction
\newcommand{\Cech}{\check{C}}              % Čech complex
\newcommand{\Hcoh}[1]{H^{#1}}             % cohomology group

% Descent
\newcommand{\descent}{\mathrm{desc}}
\newcommand{\glue}{\mathrm{glue}}
\newcommand{\overlap}[2]{#1 \cap #2}

% Semantic moves
\newcommand{\VERIFY}{\textsc{Verify}}
\newcommand{\CONSTRUCT}{\textsc{Construct}}
\newcommand{\NORMALIZE}{\textsc{Normalize}}
\newcommand{\GLUE}{\textsc{Glue}}
\newcommand{\RESTRICT}{\textsc{Restrict}}
\newcommand{\TRANSPORT}{\textsc{Transport}}
\newcommand{\REPAIR}{\textsc{Repair}}
\newcommand{\PROMOTE}{\textsc{Promote}}
\newcommand{\CHALLENGE}{\textsc{Challenge}}
\newcommand{\ESCALATE}{\textsc{Escalate}}

% Fragments
\newcommand{\QFLIA}{\texttt{QF\_LIA}}
\newcommand{\QFLRA}{\texttt{QF\_LRA}}
\newcommand{\QFBV}{\texttt{QF\_BV}}
\newcommand{\QFUF}{\texttt{QF\_UF}}

% Misc
\newcommand{\Python}{\textsc{Python}}
\newcommand{\JuGeo}{\textsc{JuGeo}}
\newcommand{\LEAN}{\textsc{Lean}\,4}
\newcommand{\Fstar}{F$^{\star}$}
\newcommand{\Coq}{\textsc{Coq}}
\newcommand{\Dafny}{\textsc{Dafny}}

% Common shortcuts
\newcommand{\ie}{i.e.\@\xspace}
\newcommand{\eg}{e.g.\@\xspace}
\newcommand{\cf}{cf.\@\xspace}
\newcommand{\wrt}{w.r.t.\@\xspace}

% Evaluation shortcuts
\newcommand{\numcases}{300}
\newcommand{\numspec}{100}
\newcommand{\numeq}{100}
\newcommand{\numbug}{100}
\newcommand{\medtime}{6.5\,ms}
\newcommand{\totaltime}{1.949\,s}


% ── Packages ────────────────────────────────────────────────────────────
\usepackage{amsmath,amssymb,amsthm,mathtools}
\usepackage{tikz-cd}
\usepackage{listings}
\usepackage{xcolor}
\usepackage{xspace}
\usepackage{booktabs}
\usepackage{multirow}
\usepackage{enumitem}
\usepackage[expansion=false]{microtype}
\usepackage{hyperref}
\usepackage{cleveref}
% \usepackage{thmtools}  % disabled: not available in this TeX installation
\usepackage{float}
\usepackage{caption}
\usepackage{subcaption}
\usepackage{tabularx}
\usepackage{stmaryrd}       % \llbracket, \rrbracket
\usepackage{bbm}            % \mathbbm{1}

% ── Page geometry ───────────────────────────────────────────────────────
\usepackage[margin=1in]{geometry}

% ── Theorem environments ────────────────────────────────────────────────
\theoremstyle{plain}
\newtheorem{theorem}{Theorem}[section]
\newtheorem{lemma}[theorem]{Lemma}
\newtheorem{proposition}[theorem]{Proposition}
\newtheorem{corollary}[theorem]{Corollary}

\theoremstyle{definition}
\newtheorem{definition}[theorem]{Definition}
\newtheorem{example}[theorem]{Example}
\newtheorem{construction}[theorem]{Construction}

\theoremstyle{remark}
\newtheorem{remark}[theorem]{Remark}
\newtheorem{notation}[theorem]{Notation}

% ── Python listings style ──────────────────────────────────────────────
\definecolor{jugeoBlue}{HTML}{2563EB}
\definecolor{jugeoGreen}{HTML}{059669}
\definecolor{jugeoRed}{HTML}{DC2626}
\definecolor{jugeoPurple}{HTML}{7C3AED}
\definecolor{jugeoGray}{HTML}{6B7280}
\definecolor{jugeoBg}{HTML}{F8FAFC}

\lstdefinestyle{jugeo-python}{
  language=Python,
  basicstyle=\ttfamily\small,
  keywordstyle=\color{jugeoBlue}\bfseries,
  stringstyle=\color{jugeoGreen},
  commentstyle=\color{jugeoGray}\itshape,
  numberstyle=\tiny\color{jugeoGray},
  backgroundcolor=\color{jugeoBg},
  numbers=left,
  numbersep=6pt,
  frame=single,
  framerule=0.4pt,
  rulecolor=\color{jugeoGray!40},
  breaklines=true,
  breakatwhitespace=true,
  showstringspaces=false,
  tabsize=4,
  xleftmargin=1.5em,
  framexleftmargin=1.5em,
  morekeywords={dataclass,frozen,field,Enum,IntEnum,auto,Any,
                Mapping,Sequence,Optional,match,case},
  literate={→}{{$\to$}}1 {←}{{$\leftarrow$}}1
           {≤}{{$\leq$}}1 {≥}{{$\geq$}}1 {≠}{{$\neq$}}1
           {∈}{{$\in$}}1 {∀}{{$\forall$}}1 {∃}{{$\exists$}}1
           {⊕}{{$\oplus$}}1 {⊖}{{$\ominus$}}1,
  escapeinside={(*@}{@*)},
}
\lstset{style=jugeo-python}

\lstdefinestyle{jugeo-output}{
  basicstyle=\ttfamily\small,
  backgroundcolor=\color{jugeoBg},
  frame=single,
  framerule=0.4pt,
  rulecolor=\color{jugeoGray!40},
  breaklines=true,
  showstringspaces=false,
  xleftmargin=1.5em,
  framexleftmargin=0.5em,
  numbers=none,
}

% ── JuGeo-specific macros ──────────────────────────────────────────────

% Judgment 8-tuple
\newcommand{\J}{\mathcal{J}}
\newcommand{\Jud}[8]{(#1,\,#2,\,#3,\,#4,\,#5,\,#6,\,#7,\,#8)}
\newcommand{\JudShort}[1]{\J_{#1}}

% Components
\newcommand{\coord}{\mathsf{c}}            % coordinate
\newcommand{\prop}{\varphi}                 % proposition
\newcommand{\carrier}{\mathcal{A}}          % carrier type
\newcommand{\evid}{\mathcal{E}}             % evidence bundle
\newcommand{\oblig}{\mathcal{O}}            % obligations
\newcommand{\obstr}{\mathcal{B}}            % obstructions
\newcommand{\trust}{\mathcal{T}}            % trust annotation
\newcommand{\prov}{\Pi}                     % provenance

% Trust algebra
\newcommand{\Talg}{\mathfrak{T}}            % trust ordered algebra
\newcommand{\Eadm}{\mathcal{E}_{\mathrm{adm}}}
\newcommand{\tleq}{\preceq}                 % trust ordering
\newcommand{\tjoin}{\oplus}                 % conservative join
\newcommand{\tatten}{\ominus}               % attenuation
\newcommand{\tpromote}{\uparrow_\pi}        % promotion
\newcommand{\tdemote}{\downarrow_\chi}       % demotion

% Trust tiers
\newcommand{\Tcontra}{\ensuremath{\mathsf{CONTRADICTED}}}
\newcommand{\Tunver}{\ensuremath{\mathsf{UNVERIFIED}}}
\newcommand{\Tcopilot}{\ensuremath{\mathsf{COPILOT}}}
\newcommand{\Toracle}{\ensuremath{\mathsf{ORACLE}}}
\newcommand{\Truntime}{\ensuremath{\mathsf{RUNTIME}}}
\newcommand{\Tsolver}{\ensuremath{\mathsf{SOLVER}}}
\newcommand{\Tproof}{\ensuremath{\mathsf{PROOF}}}

% Site and geometry
\newcommand{\Site}{\mathcal{S}}             % semantic site
\newcommand{\Cov}{\mathrm{Cov}}            % covering family
\newcommand{\Ob}{\mathrm{Ob}}              % objects
\newcommand{\Mor}{\mathrm{Mor}}            % morphisms
\newcommand{\res}{\mathrm{res}}            % restriction
\newcommand{\Cech}{\check{C}}              % Čech complex
\newcommand{\Hcoh}[1]{H^{#1}}             % cohomology group

% Descent
\newcommand{\descent}{\mathrm{desc}}
\newcommand{\glue}{\mathrm{glue}}
\newcommand{\overlap}[2]{#1 \cap #2}

% Semantic moves
\newcommand{\VERIFY}{\textsc{Verify}}
\newcommand{\CONSTRUCT}{\textsc{Construct}}
\newcommand{\NORMALIZE}{\textsc{Normalize}}
\newcommand{\GLUE}{\textsc{Glue}}
\newcommand{\RESTRICT}{\textsc{Restrict}}
\newcommand{\TRANSPORT}{\textsc{Transport}}
\newcommand{\REPAIR}{\textsc{Repair}}
\newcommand{\PROMOTE}{\textsc{Promote}}
\newcommand{\CHALLENGE}{\textsc{Challenge}}
\newcommand{\ESCALATE}{\textsc{Escalate}}

% Fragments
\newcommand{\QFLIA}{\texttt{QF\_LIA}}
\newcommand{\QFLRA}{\texttt{QF\_LRA}}
\newcommand{\QFBV}{\texttt{QF\_BV}}
\newcommand{\QFUF}{\texttt{QF\_UF}}

% Misc
\newcommand{\Python}{\textsc{Python}}
\newcommand{\JuGeo}{\textsc{JuGeo}}
\newcommand{\LEAN}{\textsc{Lean}\,4}
\newcommand{\Fstar}{F$^{\star}$}
\newcommand{\Coq}{\textsc{Coq}}
\newcommand{\Dafny}{\textsc{Dafny}}

% Common shortcuts
\newcommand{\ie}{i.e.\@\xspace}
\newcommand{\eg}{e.g.\@\xspace}
\newcommand{\cf}{cf.\@\xspace}
\newcommand{\wrt}{w.r.t.\@\xspace}

% Evaluation shortcuts
\newcommand{\numcases}{300}
\newcommand{\numspec}{100}
\newcommand{\numeq}{100}
\newcommand{\numbug}{100}
\newcommand{\medtime}{6.5\,ms}
\newcommand{\totaltime}{1.949\,s}


% ── Packages ────────────────────────────────────────────────────────────
\usepackage{amsmath,amssymb,amsthm,mathtools}
\usepackage{tikz-cd}
\usepackage{listings}
\usepackage{xcolor}
\usepackage{xspace}
\usepackage{booktabs}
\usepackage{multirow}
\usepackage{enumitem}
\usepackage[expansion=false]{microtype}
\usepackage{hyperref}
\usepackage{cleveref}
% \usepackage{thmtools}  % disabled: not available in this TeX installation
\usepackage{float}
\usepackage{caption}
\usepackage{subcaption}
\usepackage{tabularx}
\usepackage{stmaryrd}       % \llbracket, \rrbracket
\usepackage{bbm}            % \mathbbm{1}

% ── Page geometry ───────────────────────────────────────────────────────
\usepackage[margin=1in]{geometry}

% ── Theorem environments ────────────────────────────────────────────────
\theoremstyle{plain}
\newtheorem{theorem}{Theorem}[section]
\newtheorem{lemma}[theorem]{Lemma}
\newtheorem{proposition}[theorem]{Proposition}
\newtheorem{corollary}[theorem]{Corollary}

\theoremstyle{definition}
\newtheorem{definition}[theorem]{Definition}
\newtheorem{example}[theorem]{Example}
\newtheorem{construction}[theorem]{Construction}

\theoremstyle{remark}
\newtheorem{remark}[theorem]{Remark}
\newtheorem{notation}[theorem]{Notation}

% ── Python listings style ──────────────────────────────────────────────
\definecolor{jugeoBlue}{HTML}{2563EB}
\definecolor{jugeoGreen}{HTML}{059669}
\definecolor{jugeoRed}{HTML}{DC2626}
\definecolor{jugeoPurple}{HTML}{7C3AED}
\definecolor{jugeoGray}{HTML}{6B7280}
\definecolor{jugeoBg}{HTML}{F8FAFC}

\lstdefinestyle{jugeo-python}{
  language=Python,
  basicstyle=\ttfamily\small,
  keywordstyle=\color{jugeoBlue}\bfseries,
  stringstyle=\color{jugeoGreen},
  commentstyle=\color{jugeoGray}\itshape,
  numberstyle=\tiny\color{jugeoGray},
  backgroundcolor=\color{jugeoBg},
  numbers=left,
  numbersep=6pt,
  frame=single,
  framerule=0.4pt,
  rulecolor=\color{jugeoGray!40},
  breaklines=true,
  breakatwhitespace=true,
  showstringspaces=false,
  tabsize=4,
  xleftmargin=1.5em,
  framexleftmargin=1.5em,
  morekeywords={dataclass,frozen,field,Enum,IntEnum,auto,Any,
                Mapping,Sequence,Optional,match,case},
  literate={→}{{$\to$}}1 {←}{{$\leftarrow$}}1
           {≤}{{$\leq$}}1 {≥}{{$\geq$}}1 {≠}{{$\neq$}}1
           {∈}{{$\in$}}1 {∀}{{$\forall$}}1 {∃}{{$\exists$}}1
           {⊕}{{$\oplus$}}1 {⊖}{{$\ominus$}}1,
  escapeinside={(*@}{@*)},
}
\lstset{style=jugeo-python}

\lstdefinestyle{jugeo-output}{
  basicstyle=\ttfamily\small,
  backgroundcolor=\color{jugeoBg},
  frame=single,
  framerule=0.4pt,
  rulecolor=\color{jugeoGray!40},
  breaklines=true,
  showstringspaces=false,
  xleftmargin=1.5em,
  framexleftmargin=0.5em,
  numbers=none,
}

% ── JuGeo-specific macros ──────────────────────────────────────────────

% Judgment 8-tuple
\newcommand{\J}{\mathcal{J}}
\newcommand{\Jud}[8]{(#1,\,#2,\,#3,\,#4,\,#5,\,#6,\,#7,\,#8)}
\newcommand{\JudShort}[1]{\J_{#1}}

% Components
\newcommand{\coord}{\mathsf{c}}            % coordinate
\newcommand{\prop}{\varphi}                 % proposition
\newcommand{\carrier}{\mathcal{A}}          % carrier type
\newcommand{\evid}{\mathcal{E}}             % evidence bundle
\newcommand{\oblig}{\mathcal{O}}            % obligations
\newcommand{\obstr}{\mathcal{B}}            % obstructions
\newcommand{\trust}{\mathcal{T}}            % trust annotation
\newcommand{\prov}{\Pi}                     % provenance

% Trust algebra
\newcommand{\Talg}{\mathfrak{T}}            % trust ordered algebra
\newcommand{\Eadm}{\mathcal{E}_{\mathrm{adm}}}
\newcommand{\tleq}{\preceq}                 % trust ordering
\newcommand{\tjoin}{\oplus}                 % conservative join
\newcommand{\tatten}{\ominus}               % attenuation
\newcommand{\tpromote}{\uparrow_\pi}        % promotion
\newcommand{\tdemote}{\downarrow_\chi}       % demotion

% Trust tiers
\newcommand{\Tcontra}{\ensuremath{\mathsf{CONTRADICTED}}}
\newcommand{\Tunver}{\ensuremath{\mathsf{UNVERIFIED}}}
\newcommand{\Tcopilot}{\ensuremath{\mathsf{COPILOT}}}
\newcommand{\Toracle}{\ensuremath{\mathsf{ORACLE}}}
\newcommand{\Truntime}{\ensuremath{\mathsf{RUNTIME}}}
\newcommand{\Tsolver}{\ensuremath{\mathsf{SOLVER}}}
\newcommand{\Tproof}{\ensuremath{\mathsf{PROOF}}}

% Site and geometry
\newcommand{\Site}{\mathcal{S}}             % semantic site
\newcommand{\Cov}{\mathrm{Cov}}            % covering family
\newcommand{\Ob}{\mathrm{Ob}}              % objects
\newcommand{\Mor}{\mathrm{Mor}}            % morphisms
\newcommand{\res}{\mathrm{res}}            % restriction
\newcommand{\Cech}{\check{C}}              % Čech complex
\newcommand{\Hcoh}[1]{H^{#1}}             % cohomology group

% Descent
\newcommand{\descent}{\mathrm{desc}}
\newcommand{\glue}{\mathrm{glue}}
\newcommand{\overlap}[2]{#1 \cap #2}

% Semantic moves
\newcommand{\VERIFY}{\textsc{Verify}}
\newcommand{\CONSTRUCT}{\textsc{Construct}}
\newcommand{\NORMALIZE}{\textsc{Normalize}}
\newcommand{\GLUE}{\textsc{Glue}}
\newcommand{\RESTRICT}{\textsc{Restrict}}
\newcommand{\TRANSPORT}{\textsc{Transport}}
\newcommand{\REPAIR}{\textsc{Repair}}
\newcommand{\PROMOTE}{\textsc{Promote}}
\newcommand{\CHALLENGE}{\textsc{Challenge}}
\newcommand{\ESCALATE}{\textsc{Escalate}}

% Fragments
\newcommand{\QFLIA}{\texttt{QF\_LIA}}
\newcommand{\QFLRA}{\texttt{QF\_LRA}}
\newcommand{\QFBV}{\texttt{QF\_BV}}
\newcommand{\QFUF}{\texttt{QF\_UF}}

% Misc
\newcommand{\Python}{\textsc{Python}}
\newcommand{\JuGeo}{\textsc{JuGeo}}
\newcommand{\LEAN}{\textsc{Lean}\,4}
\newcommand{\Fstar}{F$^{\star}$}
\newcommand{\Coq}{\textsc{Coq}}
\newcommand{\Dafny}{\textsc{Dafny}}

% Common shortcuts
\newcommand{\ie}{i.e.\@\xspace}
\newcommand{\eg}{e.g.\@\xspace}
\newcommand{\cf}{cf.\@\xspace}
\newcommand{\wrt}{w.r.t.\@\xspace}

% Evaluation shortcuts
\newcommand{\numcases}{300}
\newcommand{\numspec}{100}
\newcommand{\numeq}{100}
\newcommand{\numbug}{100}
\newcommand{\medtime}{6.5\,ms}
\newcommand{\totaltime}{1.949\,s}

% experiment-data.tex — AUTO-GENERATED from real jugeo CLI runs
% DO NOT EDIT — regenerate with: python3 experiments/run_universal_metrics.py
% Generated from 30 programs across 6 domains

\newcommand{\expTotalPrograms}{30}
\newcommand{\expVerified}{30}
\newcommand{\expAccuracy}{100.0\%}
\newcommand{\expCoordsMin}{2}
\newcommand{\expCoordsMax}{10}
\newcommand{\expCoordsMean}{5.2}
\newcommand{\expCoordsMedian}{5.5}
\newcommand{\expPropsMin}{4}
\newcommand{\expPropsMax}{20}
\newcommand{\expPropsMean}{10.4}
\newcommand{\expPropsSum}{311}
\newcommand{\expPropsOkSum}{310}
\newcommand{\expTimeMin}{3.506\,s}
\newcommand{\expTimeMax}{6.714\,s}
\newcommand{\expTimeMean}{4.64\,s}
\newcommand{\expTimeMedian}{4.508\,s}
\newcommand{\expTimeTotal}{139.204\,s}
\newcommand{\expObstructionTotal}{0}
\newcommand{\expHOneZero}{30}
\newcommand{\expDomainCount}{6}
\newcommand{\expTrustCOPILOTSUGGESTED}{30}
\newcommand{\expDomainDsCount}{5}
\newcommand{\expDomainDsVerified}{5}
\newcommand{\expDomainDsAvgCoords}{7.6}
\newcommand{\expDomainDsAvgProps}{15.2}
\newcommand{\expDomainMathCount}{5}
\newcommand{\expDomainMathVerified}{5}
\newcommand{\expDomainMathAvgCoords}{4.8}
\newcommand{\expDomainMathAvgProps}{9.4}
\newcommand{\expDomainSmCount}{5}
\newcommand{\expDomainSmVerified}{5}
\newcommand{\expDomainSmAvgCoords}{7.6}
\newcommand{\expDomainSmAvgProps}{15.2}
\newcommand{\expDomainSortCount}{5}
\newcommand{\expDomainSortVerified}{5}
\newcommand{\expDomainSortAvgCoords}{2.4}
\newcommand{\expDomainSortAvgProps}{4.8}
\newcommand{\expDomainStrCount}{5}
\newcommand{\expDomainStrVerified}{5}
\newcommand{\expDomainStrAvgCoords}{3.2}
\newcommand{\expDomainStrAvgProps}{6.4}
\newcommand{\expDomainWebCount}{5}
\newcommand{\expDomainWebVerified}{5}
\newcommand{\expDomainWebAvgCoords}{5.6}
\newcommand{\expDomainWebAvgProps}{11.2}

% subsystem-data.tex — AUTO-GENERATED from real jugeo subsystem experiments
% DO NOT EDIT — regenerate with: python3 experiments/run_subsystem_experiments.py

\newcommand{\subCliTotal}{10}
\newcommand{\subCliVerified}{9}
\newcommand{\subCliAccuracy}{90.0\%}
\newcommand{\subCliMeanTime}{4.132\,s}
\newcommand{\subCliMinTime}{3.472\,s}
\newcommand{\subCliMaxTime}{5.372\,s}
\newcommand{\subCliTotalCoords}{54}
\newcommand{\subCliTotalProps}{107}
\newcommand{\subCliTotalPropsOk}{106}
\newcommand{\subSiteCalls}{90}
\newcommand{\subSiteSuccessful}{69}
\newcommand{\subSiteSuccessRate}{76.7\%}
\newcommand{\subSiteMeanTime}{0.0001\,s}
\newcommand{\subSiteTotalTime}{0.005\,s}
\newcommand{\subSpecificationsatisfactionSmallTime}{0.0003\,s}
\newcommand{\subSpecificationsatisfactionMediumTime}{0.0001\,s}
\newcommand{\subSpecificationsatisfactionLargeTime}{0.0001\,s}
\newcommand{\subInhabitantfleetSmallTime}{0.0\,s}
\newcommand{\subInhabitantfleetMediumTime}{0.0\,s}
\newcommand{\subInhabitantfleetLargeTime}{0.0\,s}
\newcommand{\subTheoremecologySmallTime}{0.0\,s}
\newcommand{\subTheoremecologyMediumTime}{0.0\,s}
\newcommand{\subTheoremecologyLargeTime}{0.0\,s}
\newcommand{\subStatespaceexplorationSmallTime}{0.0\,s}
\newcommand{\subStatespaceexplorationMediumTime}{0.0\,s}
\newcommand{\subStatespaceexplorationLargeTime}{0.0\,s}
\newcommand{\subRepairsemanticsSmallTime}{0.0\,s}
\newcommand{\subRepairsemanticsMediumTime}{0.0\,s}
\newcommand{\subRepairsemanticsLargeTime}{0.0\,s}
\newcommand{\subReplaygluingSmallTime}{0.0\,s}
\newcommand{\subReplaygluingMediumTime}{0.0\,s}
\newcommand{\subReplaygluingLargeTime}{0.0\,s}
\newcommand{\subSemanticfuturesSmallTime}{0.0\,s}
\newcommand{\subSemanticfuturesMediumTime}{0.0\,s}
\newcommand{\subSemanticfuturesLargeTime}{0.0\,s}
\newcommand{\subHypercovertreatySmallTime}{0.0\,s}
\newcommand{\subHypercovertreatyMediumTime}{0.0\,s}
\newcommand{\subHypercovertreatyLargeTime}{0.0\,s}
\newcommand{\subEncodeforsolverSmallTime}{0.0026\,s}
\newcommand{\subEncodeforsolverMediumTime}{0.0002\,s}
\newcommand{\subEncodeforsolverLargeTime}{0.0001\,s}
\newcommand{\subInterfaceroutingSmallTime}{0.0\,s}
\newcommand{\subInterfaceroutingMediumTime}{0.0\,s}
\newcommand{\subInterfaceroutingLargeTime}{0.0\,s}
\newcommand{\subOrchestrateverificationSmallTime}{0.0002\,s}
\newcommand{\subOrchestrateverificationMediumTime}{0.0\,s}
\newcommand{\subOrchestrateverificationLargeTime}{0.0\,s}
\newcommand{\subRunfulldescentSmallTime}{0.0\,s}
\newcommand{\subRunfulldescentMediumTime}{0.0\,s}
\newcommand{\subRunfulldescentLargeTime}{0.0\,s}
\newcommand{\subKernellifecycleSmallTime}{0.0\,s}
\newcommand{\subKernellifecycleMediumTime}{0.0\,s}
\newcommand{\subKernellifecycleLargeTime}{0.0\,s}
\newcommand{\subDiscoverypipelineSmallTime}{0.0\,s}
\newcommand{\subDiscoverypipelineMediumTime}{0.0\,s}
\newcommand{\subDiscoverypipelineLargeTime}{0.0\,s}
\newcommand{\subAnalogytransportSmallTime}{0.0\,s}
\newcommand{\subAnalogytransportMediumTime}{0.0\,s}
\newcommand{\subAnalogytransportLargeTime}{0.0\,s}
\newcommand{\subFormalcoresiteSmallTime}{0.0001\,s}
\newcommand{\subFormalcoresiteMediumTime}{0.0\,s}
\newcommand{\subFormalcoresiteLargeTime}{0.0\,s}
\newcommand{\subProblematlasSmallTime}{0.0\,s}
\newcommand{\subProblematlasMediumTime}{0.0\,s}
\newcommand{\subProblematlasLargeTime}{0.0\,s}
\newcommand{\subPublicalignmentSmallTime}{0.0\,s}
\newcommand{\subPublicalignmentMediumTime}{0.0\,s}
\newcommand{\subPublicalignmentLargeTime}{0.0\,s}
\newcommand{\subRegimebootstrappingSmallTime}{0.0\,s}
\newcommand{\subRegimebootstrappingMediumTime}{0.0\,s}
\newcommand{\subRegimebootstrappingLargeTime}{0.0\,s}
\newcommand{\subRelationalrefinementSmallTime}{0.0\,s}
\newcommand{\subRelationalrefinementMediumTime}{0.0\,s}
\newcommand{\subRelationalrefinementLargeTime}{0.0\,s}
\newcommand{\subSemanticclosureSmallTime}{0.0\,s}
\newcommand{\subSemanticclosureMediumTime}{0.0\,s}
\newcommand{\subSemanticclosureLargeTime}{0.0\,s}
\newcommand{\subTheoremeconomicsSmallTime}{0.0001\,s}
\newcommand{\subTheoremeconomicsMediumTime}{0.0\,s}
\newcommand{\subTheoremeconomicsLargeTime}{0.0\,s}
\newcommand{\subBenchmarksuiteSmallTime}{0.0\,s}
\newcommand{\subBenchmarksuiteMediumTime}{0.0\,s}
\newcommand{\subBenchmarksuiteLargeTime}{0.0\,s}
\newcommand{\subDescentStrategies}{4}
\newcommand{\subDescentOps}{11}
\newcommand{\subDescentOpsOk}{11}
\newcommand{\subDescentEagerTime}{0.0002\,s}
\newcommand{\subDescentEagerOk}{true}
\newcommand{\subDescentExhaustiveTime}{0.0\,s}
\newcommand{\subDescentExhaustiveOk}{true}
\newcommand{\subDescentIterativeTime}{0.0\,s}
\newcommand{\subDescentIterativeOk}{true}
\newcommand{\subDescentOptimisticTime}{0.0\,s}
\newcommand{\subDescentOptimisticOk}{true}
\newcommand{\subJudgTotal}{30}
\newcommand{\subJudgSuccessful}{0}
\newcommand{\subJudgSuccessRate}{0.0\%}
\newcommand{\subJudgMeanTimeUs}{7.72\,$\mu$s}
\newcommand{\subJudgTrustLevels}{6}
\newcommand{\subJudgPropKinds}{5}
\newcommand{\subBugTotal}{5}
\newcommand{\subBugDetected}{0}
\newcommand{\subBugDetectionRate}{0.0\%}
\newcommand{\subBugFalsePositiveRate}{0.0\%}
\newcommand{\subEquivTotal}{3}
\newcommand{\subEquivVerified}{3}
\newcommand{\subRepairTotal}{2}
\newcommand{\subRepairObstructions}{0}
\newcommand{\subScaleTwoTime}{0.0001\,s}
\newcommand{\subScaleFourTime}{0.0\,s}
\newcommand{\subScaleEightTime}{0.0\,s}
\newcommand{\subScaleTwelveTime}{0.0\,s}
\newcommand{\subScaleSixteenTime}{0.0\,s}
\newcommand{\subScaleTwentyTime}{0.0\,s}
\newcommand{\subTotalExperimentTime}{95.6\,s}

% comprehensive-data.tex — AUTO-GENERATED from real jugeo experiments
% DO NOT EDIT — regenerate with: python3 experiments/run_comprehensive_experiments.py
% Generated: 2026-03-23 09:19:06

% --- Size scaling metrics ---
\newcommand{\compTotalPrograms}{18}
\newcommand{\compTotalVerified}{18}
\newcommand{\compOverallAccuracy}{100.0\%}
\newcommand{\compTinyCount}{5}
\newcommand{\compTinyVerified}{5}
\newcommand{\compTinyMeanCoords}{3}
\newcommand{\compTinyMeanProps}{6}
\newcommand{\compTinyMeanTime}{4.95\,s}
\newcommand{\compTinyMeanLines}{5}
\newcommand{\compTinyMinTime}{4.45\,s}
\newcommand{\compTinyMaxTime}{5.51\,s}
\newcommand{\compSmallCount}{5}
\newcommand{\compSmallVerified}{5}
\newcommand{\compSmallMeanCoords}{3.6}
\newcommand{\compSmallMeanProps}{7.2}
\newcommand{\compSmallMeanTime}{4.85\,s}
\newcommand{\compSmallMeanLines}{16.0}
\newcommand{\compSmallMinTime}{4.30\,s}
\newcommand{\compSmallMaxTime}{5.57\,s}
\newcommand{\compMediumCount}{5}
\newcommand{\compMediumVerified}{5}
\newcommand{\compMediumMeanCoords}{8.6}
\newcommand{\compMediumMeanProps}{17.2}
\newcommand{\compMediumMeanTime}{4.47\,s}
\newcommand{\compMediumMeanLines}{45.0}
\newcommand{\compMediumMinTime}{4.26\,s}
\newcommand{\compMediumMaxTime}{4.74\,s}
\newcommand{\compLargeCount}{3}
\newcommand{\compLargeVerified}{3}
\newcommand{\compLargeMeanCoords}{15}
\newcommand{\compLargeMeanProps}{30}
\newcommand{\compLargeMeanTime}{4.31\,s}
\newcommand{\compLargeMeanLines}{79.0}
\newcommand{\compLargeMinTime}{4.27\,s}
\newcommand{\compLargeMaxTime}{4.38\,s}
% --- Aggregate timing ---
\newcommand{\compTimeMean}{4.68\,s}
\newcommand{\compTimeMedian}{4.50\,s}
\newcommand{\compTimeMin}{4.265\,s}
\newcommand{\compTimeMax}{5.571\,s}
\newcommand{\compTimeTotal}{84.3\,s}
\newcommand{\compTimeStdev}{0.45\,s}
% --- Aggregate structure ---
\newcommand{\compCoordsMean}{6.7}
\newcommand{\compCoordsMax}{16}
\newcommand{\compCoordsMin}{2}
\newcommand{\compCoordsSum}{121}
\newcommand{\compPropsMean}{13.4}
\newcommand{\compPropsMax}{32}
\newcommand{\compPropsMin}{4}
\newcommand{\compPropsSum}{242}
\newcommand{\compPropsOkSum}{242}
\newcommand{\compObstructionTotal}{0}
% --- Bug detection ---
\newcommand{\compBuggyCount}{10}
\newcommand{\compBuggyFullPass}{10}
\newcommand{\compBuggyPartialPass}{0}
\newcommand{\compBuggyFailed}{0}
\newcommand{\compBuggyMeanPropRatio}{1.00}
\newcommand{\compCorrectMeanPropRatio}{1.00}
\newcommand{\compBuggyMeanTime}{5.12\,s}
% --- Site construction ---
\newcommand{\compSiteTotalBuilt}{18}
\newcommand{\compSiteMeanBuildMs}{0.05}
\newcommand{\compSiteMaxBuildMs}{0.14}
\newcommand{\compSiteMeanCoords}{6.5}
\newcommand{\compSiteMeanMorphisms}{14.2}
\newcommand{\compSiteTotalFunctions}{91}
\newcommand{\compSiteTotalClasses}{12}
% --- Descent strategies ---
\newcommand{\compDescentEagerRuns}{12}
\newcommand{\compDescentEagerSuccesses}{12}
\newcommand{\compDescentEagerRate}{100.0\%}
\newcommand{\compDescentEagerMeanMs}{0.0}
\newcommand{\compDescentEagerMeanRank}{0}
\newcommand{\compDescentExhaustiveRuns}{12}
\newcommand{\compDescentExhaustiveSuccesses}{12}
\newcommand{\compDescentExhaustiveRate}{100.0\%}
\newcommand{\compDescentExhaustiveMeanMs}{0.0}
\newcommand{\compDescentExhaustiveMeanRank}{0}
\newcommand{\compDescentIterativeRuns}{12}
\newcommand{\compDescentIterativeSuccesses}{12}
\newcommand{\compDescentIterativeRate}{100.0\%}
\newcommand{\compDescentIterativeMeanMs}{0.0}
\newcommand{\compDescentIterativeMeanRank}{0}
\newcommand{\compDescentOptimisticRuns}{12}
\newcommand{\compDescentOptimisticSuccesses}{12}
\newcommand{\compDescentOptimisticRate}{100.0\%}
\newcommand{\compDescentOptimisticMeanMs}{0.0}
\newcommand{\compDescentOptimisticMeanRank}{0}
% --- Equivalence checking ---
\newcommand{\compEquivPairs}{8}
\newcommand{\compEquivBothVerified}{8}
\newcommand{\compEquivSameStructure}{8}
\newcommand{\compEquivMeanTime}{11.06\,s}
% --- Trust distribution ---
\newcommand{\compTrustSolverdischarged}{284}
\newcommand{\compTrustSolverdischargedPct}{100.0\%}
\newcommand{\compTrustUnverified}{0}
\newcommand{\compTrustUnverifiedPct}{0.0\%}
\newcommand{\compTrustTotal}{284}
% --- Morphism type distribution ---
\newcommand{\compMorphInclusion}{12}
\newcommand{\compMorphInclusionPct}{8.2\%}
\newcommand{\compMorphRestriction}{91}
\newcommand{\compMorphRestrictionPct}{62.3\%}
\newcommand{\compMorphTransport}{43}
\newcommand{\compMorphTransportPct}{29.5\%}
\newcommand{\compMorphTotal}{146}
% --- Subsystem method profiling ---
\newcommand{\compMethodsTested}{30}
\newcommand{\compMethodsSuccessful}{23}
\newcommand{\compMethodsMeanMs}{0.02}
\newcommand{\compProfAnalogyTransportMeanMs}{0.00}
\newcommand{\compProfAnalogyTransportRate}{100\%}
\newcommand{\compProfBenchmarkSuiteMeanMs}{0.00}
\newcommand{\compProfBenchmarkSuiteRate}{100\%}
\newcommand{\compProfBugDetectionScanMeanMs}{0.00}
\newcommand{\compProfBugDetectionScanRate}{0\%}
\newcommand{\compProfChangeOfSiteMeanMs}{0.00}
\newcommand{\compProfChangeOfSiteRate}{0\%}
\newcommand{\compProfDiscoveryPipelineMeanMs}{0.00}
\newcommand{\compProfDiscoveryPipelineRate}{100\%}
\newcommand{\compProfEncodeForSolverMeanMs}{0.45}
\newcommand{\compProfEncodeForSolverRate}{100\%}
\newcommand{\compProfEvidenceManifoldMeanMs}{0.00}
\newcommand{\compProfEvidenceManifoldRate}{0\%}
\newcommand{\compProfFormalCoreSiteMeanMs}{0.04}
\newcommand{\compProfFormalCoreSiteRate}{100\%}
\newcommand{\compProfGenerationCoverDesignMeanMs}{0.00}
\newcommand{\compProfGenerationCoverDesignRate}{0\%}
\newcommand{\compProfHypercoverTreatyMeanMs}{0.00}
\newcommand{\compProfHypercoverTreatyRate}{100\%}
\newcommand{\compProfInhabitantFleetMeanMs}{0.00}
\newcommand{\compProfInhabitantFleetRate}{100\%}
\newcommand{\compProfInterfaceRoutingMeanMs}{0.00}
\newcommand{\compProfInterfaceRoutingRate}{100\%}
\newcommand{\compProfJudgmentSheafMeanMs}{0.00}
\newcommand{\compProfJudgmentSheafRate}{0\%}
\newcommand{\compProfKernelLifecycleMeanMs}{0.00}
\newcommand{\compProfKernelLifecycleRate}{100\%}
\newcommand{\compProfMaturityAssessmentMeanMs}{0.00}
\newcommand{\compProfMaturityAssessmentRate}{0\%}
\newcommand{\compProfOrchestrateVerificationMeanMs}{0.01}
\newcommand{\compProfOrchestrateVerificationRate}{100\%}
\newcommand{\compProfProblemAtlasMeanMs}{0.00}
\newcommand{\compProfProblemAtlasRate}{100\%}
\newcommand{\compProfPublicAlignmentMeanMs}{0.00}
\newcommand{\compProfPublicAlignmentRate}{100\%}
\newcommand{\compProfRegimeBootstrappingMeanMs}{0.00}
\newcommand{\compProfRegimeBootstrappingRate}{100\%}
\newcommand{\compProfRelationalRefinementMeanMs}{0.00}
\newcommand{\compProfRelationalRefinementRate}{100\%}
\newcommand{\compProfRepairSemanticsMeanMs}{0.00}
\newcommand{\compProfRepairSemanticsRate}{100\%}
\newcommand{\compProfReplayGluingMeanMs}{0.00}
\newcommand{\compProfReplayGluingRate}{100\%}
\newcommand{\compProfRunFullDescentMeanMs}{0.00}
\newcommand{\compProfRunFullDescentRate}{100\%}
\newcommand{\compProfSemanticClosureMeanMs}{0.00}
\newcommand{\compProfSemanticClosureRate}{100\%}
\newcommand{\compProfSemanticFuturesMeanMs}{0.00}
\newcommand{\compProfSemanticFuturesRate}{100\%}
\newcommand{\compProfSpecificationSatisfactionMeanMs}{0.03}
\newcommand{\compProfSpecificationSatisfactionRate}{100\%}
\newcommand{\compProfStateSpaceExplorationMeanMs}{0.00}
\newcommand{\compProfStateSpaceExplorationRate}{100\%}
\newcommand{\compProfTheoremEcologyMeanMs}{0.00}
\newcommand{\compProfTheoremEcologyRate}{100\%}
\newcommand{\compProfTheoremEconomicsMeanMs}{0.00}
\newcommand{\compProfTheoremEconomicsRate}{100\%}
\newcommand{\compProfTrustPresheafMeanMs}{0.00}
\newcommand{\compProfTrustPresheafRate}{0\%}

% data-paper26.tex -- AUTO-GENERATED by exp26_import_graph.py
% DO NOT EDIT -- regenerate with: python3 experiments/exp26_import_graph.py

% --- Overall site structure ---
\newcommand{\ppTwentysixTotalPrograms}{10}
\newcommand{\ppTwentysixTotalCoords}{77}
\newcommand{\ppTwentysixTotalMorphisms}{114}
\newcommand{\ppTwentysixMeanCoords}{7.7}
\newcommand{\ppTwentysixMeanMorphisms}{11.4}
\newcommand{\ppTwentysixMeanDepth}{3.0}
\newcommand{\ppTwentysixMaxDepth}{4}

% --- Verification results ---
\newcommand{\ppTwentysixVerifiedCount}{10}
\newcommand{\ppTwentysixTotalObstructions}{0}
\newcommand{\ppTwentysixTotalOverlap}{293}

% --- Unordered strategy ---
\newcommand{\ppTwentysixUnordReverif}{38}
\newcommand{\ppTwentysixUnordMeanReverif}{3.8}
\newcommand{\ppTwentysixUnordTime}{61718.20\,\text{ms}}
\newcommand{\ppTwentysixUnordMeanTime}{6171.82\,\text{ms}}

% --- Topological strategy ---
\newcommand{\ppTwentysixTopoReverif}{0}
\newcommand{\ppTwentysixTopoTime}{59855.19\,\text{ms}}
\newcommand{\ppTwentysixTopoMeanTime}{5985.52\,\text{ms}}

% --- Comparison ---
\newcommand{\ppTwentysixSpeedup}{1.03$\times$}

% --- Per-program results ---
\newcommand{\ppTwentysixLinearChainCoords}{6}
\newcommand{\ppTwentysixLinearChainMorphisms}{9}
\newcommand{\ppTwentysixLinearChainDepth}{4}
\newcommand{\ppTwentysixLinearChainUnordTime}{5071.08\,\text{ms}}
\newcommand{\ppTwentysixLinearChainTopoTime}{5299.83\,\text{ms}}
\newcommand{\ppTwentysixLinearChainReverifs}{3}
\newcommand{\ppTwentysixFanOutCoords}{7}
\newcommand{\ppTwentysixFanOutMorphisms}{14}
\newcommand{\ppTwentysixFanOutDepth}{2}
\newcommand{\ppTwentysixFanOutUnordTime}{5065.43\,\text{ms}}
\newcommand{\ppTwentysixFanOutTopoTime}{5709.06\,\text{ms}}
\newcommand{\ppTwentysixFanOutReverifs}{7}
\newcommand{\ppTwentysixFanInCoords}{6}
\newcommand{\ppTwentysixFanInMorphisms}{9}
\newcommand{\ppTwentysixFanInDepth}{2}
\newcommand{\ppTwentysixFanInUnordTime}{6068.00\,\text{ms}}
\newcommand{\ppTwentysixFanInTopoTime}{5747.03\,\text{ms}}
\newcommand{\ppTwentysixFanInReverifs}{3}
\newcommand{\ppTwentysixDiamondDepsCoords}{6}
\newcommand{\ppTwentysixDiamondDepsMorphisms}{11}
\newcommand{\ppTwentysixDiamondDepsDepth}{3}
\newcommand{\ppTwentysixDiamondDepsUnordTime}{6047.32\,\text{ms}}
\newcommand{\ppTwentysixDiamondDepsTopoTime}{5142.15\,\text{ms}}
\newcommand{\ppTwentysixDiamondDepsReverifs}{5}
\newcommand{\ppTwentysixDeepDagCoords}{8}
\newcommand{\ppTwentysixDeepDagMorphisms}{16}
\newcommand{\ppTwentysixDeepDagDepth}{4}
\newcommand{\ppTwentysixDeepDagUnordTime}{5325.92\,\text{ms}}
\newcommand{\ppTwentysixDeepDagTopoTime}{4702.95\,\text{ms}}
\newcommand{\ppTwentysixDeepDagReverifs}{8}
\newcommand{\ppTwentysixClassCallGraphCoords}{11}
\newcommand{\ppTwentysixClassCallGraphMorphisms}{10}
\newcommand{\ppTwentysixClassCallGraphDepth}{3}
\newcommand{\ppTwentysixClassCallGraphUnordTime}{6188.93\,\text{ms}}
\newcommand{\ppTwentysixClassCallGraphTopoTime}{6003.42\,\text{ms}}
\newcommand{\ppTwentysixClassCallGraphReverifs}{0}
\newcommand{\ppTwentysixMutualHelpersCoords}{5}
\newcommand{\ppTwentysixMutualHelpersMorphisms}{8}
\newcommand{\ppTwentysixMutualHelpersDepth}{2}
\newcommand{\ppTwentysixMutualHelpersUnordTime}{6125.49\,\text{ms}}
\newcommand{\ppTwentysixMutualHelpersTopoTime}{5298.73\,\text{ms}}
\newcommand{\ppTwentysixMutualHelpersReverifs}{3}
\newcommand{\ppTwentysixUtilityLayersCoords}{6}
\newcommand{\ppTwentysixUtilityLayersMorphisms}{9}
\newcommand{\ppTwentysixUtilityLayersDepth}{4}
\newcommand{\ppTwentysixUtilityLayersUnordTime}{5095.27\,\text{ms}}
\newcommand{\ppTwentysixUtilityLayersTopoTime}{4644.46\,\text{ms}}
\newcommand{\ppTwentysixUtilityLayersReverifs}{3}
\newcommand{\ppTwentysixRegistryPatternCoords}{8}
\newcommand{\ppTwentysixRegistryPatternMorphisms}{11}
\newcommand{\ppTwentysixRegistryPatternDepth}{2}
\newcommand{\ppTwentysixRegistryPatternUnordTime}{7151.21\,\text{ms}}
\newcommand{\ppTwentysixRegistryPatternTopoTime}{7740.76\,\text{ms}}
\newcommand{\ppTwentysixRegistryPatternReverifs}{3}
\newcommand{\ppTwentysixMultiModuleSimCoords}{14}
\newcommand{\ppTwentysixMultiModuleSimMorphisms}{17}
\newcommand{\ppTwentysixMultiModuleSimDepth}{4}
\newcommand{\ppTwentysixMultiModuleSimUnordTime}{9579.56\,\text{ms}}
\newcommand{\ppTwentysixMultiModuleSimTopoTime}{9566.81\,\text{ms}}
\newcommand{\ppTwentysixMultiModuleSimReverifs}{3}


% ── Paper-specific macros ──────────────────────────────────────────────
\newcommand{\IG}{\mathcal{G}}              % import graph
\newcommand{\DA}{\textsc{DepAnalyzer}}     % DependencyAnalyzer
\newcommand{\Mod}{\mathsf{Mod}}            % a module identifier
\newcommand{\MdCat}{\mathbf{Mod}}          % module category
\newcommand{\Dep}{\mathrm{Dep}}            % direct dependency set
\newcommand{\SCCP}{\mathsf{SCC}}           % strongly connected component
\newcommand{\TopoOrd}{\tau}                % topological order symbol
\newcommand{\VerOrd}{\leq_{\tau}}          % verification order relation
\newcommand{\ImEdge}{\mathrm{imp}}         % import edge relation
\newcommand{\CyclObs}{\mathcal{C}}         % cycle obstruction
\newcommand{\LocalVerif}{\mathrm{LV}}      % local verifier
\newcommand{\GlobVerif}{\mathrm{GV}}       % global verification
\newcommand{\CondGraph}{\widehat{\IG}}     % condensation graph
\newcommand{\VerifCtx}{\Gamma}             % verification context
\newcommand{\ModSite}{\mathcal{S}_{\IG}}   % module site

\newcommand{\NN}{\mathbb{N}}
\newcommand{\ZZ}{\mathbb{Z}}
\newcommand{\tgeq}{\succeq}              % trust greater-or-equal

% ──────────────────────────────────────────────────────────────────────

\title{Import Graph Analysis:\\
       Dependency-Aware Site Construction\\[4pt]
       {\large Judgment Geometry Series, Paper~26}}
\author{Halley Young\\[2pt] Microsoft Research}
\date{2025}

\begin{document}
\maketitle

\begin{abstract}
The import graph of a Python project encodes the dependency structure
between modules: each module is a node, and each import statement is a
directed edge from importer to importee.  \JuGeo{} lifts this structure
into the semantic site framework by treating modules as coordinate objects
and imports as restriction morphisms.  We introduce \DA{}, which builds
the import graph from Python AST traversal, detects strongly-connected
components (cycles), and computes a topological verification order via
Kahn's algorithm.  The resulting \emph{module site} $\ModSite$ supports
bottom-up verification: a module is processed only after all its
transitive dependencies carry verified judgments.  The central result,
the \emph{Compositionality Theorem}, states that bottom-up verification
in topological order is sound and complete: the global verification
succeeds if and only if every local verification succeeds.  Cyclic
imports—permitted by Python but dangerous for verification—are handled
by SCC condensation: each cycle produces an obstruction record and is
resolved by a bundled joint judgment for the entire component.  We give
a full \LEAN{} formalization (\texttt{Paper26\_ImportGraph.lean}) with
the compositionality theorem proved without any \texttt{sorry} or
axioms beyond Lean's kernel.
\end{abstract}


% ══════════════════════════════════════════════════════════════════════
\section{Introduction}
\label{sec:intro}
% ══════════════════════════════════════════════════════════════════════

A Python project is not a monolith.  It is a collection of modules,
each potentially importing from others, forming a directed dependency
graph.  This graph imposes a natural partial order on verification: to
verify module $M$, one should first understand—and ideally verify—every
module that $M$ imports.  Ignoring this structure forces verification
tools to re-derive dependencies each time they process a module, missing
the opportunity to cache and reuse earlier proofs.

\JuGeo{} exploits the import graph explicitly.  The \emph{import graph}
$\IG = (V, E)$ has vertex set $V$ (the modules of the project) and edge
set $E$ (the import statements).  A directed edge $(M, N) \in E$ means
\emph{``module $M$ imports from module $N$,''} making $N$ a dependency
of $M$.  This graph is not merely auxiliary metadata; it is the
combinatorial skeleton of the verification site $\ModSite$ built by the
JuGeo runtime.

\paragraph{The central insight.}
\begin{quote}
  \itshape
  The import graph $\IG$ of a Python project is the combinatorial
  shadow of the semantic site $\ModSite$.  Modules are coordinate
  objects, import edges are restriction morphisms, and the topological
  order of $\IG$ prescribes the verification schedule.
\end{quote}

This paper develops this idea precisely.  Section~\ref{sec:construction}
describes how \JuGeo{} constructs $\IG$ from Python AST nodes.
Section~\ref{sec:analysis} presents \DA{}, which runs Tarjan's SCC
algorithm, computes the condensation $\CondGraph$, and emits a
topological order $\TopoOrd$ via Kahn's BFS.
Section~\ref{sec:site} defines the Grothendieck site $\ModSite$ and
relates its sheaf condition to compositional verification.
Section~\ref{sec:ordering} formalises the bottom-up traversal.
Section~\ref{sec:cycles} handles cyclic imports.
Section~\ref{sec:theorem} states and proves the Compositionality Theorem
and presents its \LEAN{} formalization.
Section~\ref{sec:experiments} reports benchmarks.

\paragraph{Contributions.}
\begin{enumerate}[label=(\roman*)]
  \item \textbf{Import graph construction}
        (\S\ref{sec:construction}).
        \texttt{ImportGraphBuilder} walks a Python project directory,
        parses each \texttt{.py} file with \texttt{ast}, and populates
        $\IG = (V, E)$ incrementally.

  \item \textbf{Dependency analysis}
        (\S\ref{sec:analysis}).
        \DA{} detects SCCs in $O(|V|+|E|)$ using Tarjan's algorithm,
        condenses $\IG$ to the DAG $\CondGraph$, and computes
        a topological order $\TopoOrd$.

  \item \textbf{Module site construction}
        (\S\ref{sec:site}).
        We define $\ModSite$ as a Grothendieck site whose covering
        sieves arise from import edges and show that the verification
        sheaf satisfies descent.

  \item \textbf{Verification ordering}
        (\S\ref{sec:ordering}).
        The topological order $\TopoOrd$ ensures every module is
        processed after its dependencies carry verified judgments.

  \item \textbf{Cycle handling}
        (\S\ref{sec:cycles}).
        Cyclic SCCs produce obstruction records at trust level
        $\Tsolver$; the SCC is verified as an atomic bundle judgment.

  \item \textbf{Compositionality theorem}
        (\S\ref{sec:theorem}).
        Topological-order verification is sound and complete; the
        theorem is mechanized in \LEAN{} without \texttt{sorry}.
\end{enumerate}

\paragraph{Relation to prior work.}
Papers~1--3 develop semantic sites and descent obstructions
\cite{paper01,paper02,paper03}.  Paper~4 supplies the trust algebra
\cite{paper04}.  Paper~5 handles SMT dispatch for cycle queries
\cite{paper05}.  Paper~6 defines semantic moves including \GLUE{} and
\REPAIR{} \cite{paper06}.  Paper~7 treats Python-specific effects
\cite{paper07}.  The present paper is the first in the series to study
how the \emph{module-level} structure of a Python project shapes its
verification site.  Related tools—\Dafny{}'s dependency tracker
\cite{dafny}, Rust's crate graph, Haskell's \texttt{ghc --make}—all
implicitly perform bottom-up verification; \JuGeo{} makes this structure
explicit, categorical, and proof-carrying.

% ══════════════════════════════════════════════════════════════════════
\section{Import Graph Construction}
\label{sec:construction}
% ══════════════════════════════════════════════════════════════════════

\subsection{Python import statements and the AST}

Python's standard library provides the \texttt{ast} module for parsing
source files into abstract syntax trees.  Import statements appear as
two AST node types:
\begin{itemize}
  \item \texttt{ast.Import}: absolute imports such as
        \verb|import os| and \verb|import numpy as np|.
  \item \texttt{ast.ImportFrom}: relative and absolute from-imports
        such as \verb|from pathlib import Path| and
        \verb|from . import utils|.
\end{itemize}

\texttt{ImportGraphBuilder} traverses the project root with
\texttt{os.walk}, collects every \texttt{.py} file, and calls
\texttt{ast.parse} on each.  It then walks the resulting AST,
visiting both import node types to extract edges.  If parsing fails
(malformed source), the file is skipped and a
$\trust = \Tunver$ stub node is inserted.

\begin{lstlisting}[style=jugeo-python,
  caption={Core of \texttt{ImportGraphBuilder.build}.},
  label={lst:build}]
def build(self, root: Path) -> ImportGraph:
    nodes: dict[str, ImportNode] = {}
    edges: list[ImportEdge] = []
    for dirpath, _, filenames in os.walk(root):
        for fname in filenames:
            if not fname.endswith(".py"):
                continue
            path = Path(dirpath) / fname
            mod  = self._path_to_module(path, root)
            nodes[mod] = ImportNode(
                name=mod, path=path, trust=TrustLevel.UNVERIFIED)
            try:
                tree = ast.parse(path.read_text())
            except SyntaxError:
                continue
            for node in ast.walk(tree):
                if isinstance(node, ast.Import):
                    for alias in node.names:
                        edges.append(
                            ImportEdge(src=mod, dst=alias.name))
                elif isinstance(node, ast.ImportFrom):
                    base = node.module or ""
                    edges.append(ImportEdge(src=mod, dst=base))
    return ImportGraph(nodes=nodes, edges=edges)
\end{lstlisting}

\subsection{Formal definition}

\begin{definition}[Import graph]
\label{def:import-graph}
An \emph{import graph} is a directed graph $\IG = (V, E)$ where
\begin{itemize}
  \item $V$ is a finite set of \emph{module identifiers}
        (dot-separated names such as \texttt{jugeo.runtime.site});
  \item $E \subseteq V \times V$ is the \emph{import edge} relation;
        $(M, N) \in E$ iff module $M$ contains an import statement
        referencing module $N$.
\end{itemize}
The \emph{direct dependency set} of $M$ is
$\Dep(M) \coloneqq \{N \mid (M, N) \in E\}$
and the \emph{transitive dependency set} is the reflexive-transitive
closure $\Dep^*(M)$.
\end{definition}

\begin{definition}[Well-formed import graph]
\label{def:well-formed}
$\IG$ is \emph{well-formed} if both endpoints of every edge are in $V$:
\[
  \forall (M, N) \in E,\quad M \in V \text{ and } N \in V.
\]
\end{definition}

Real projects may import third-party packages not present in the
project tree.  \texttt{ImportGraphBuilder} inserts stub nodes for
unresolved names with $\trust = \Tunver$, ensuring well-formedness.

\begin{definition}[Acyclic import graph]
\label{def:acyclic}
$\IG$ is \emph{acyclic} if the relation $E$ contains no directed cycle.
Equivalently, there exists a total strict order $<_{\TopoOrd}$ on $V$
such that $(M, N) \in E \Rightarrow N <_{\TopoOrd} M$.
\end{definition}

Real Python projects often contain cycles (see \S\ref{sec:cycles}).
\DA{} handles these via SCC condensation before computing $\TopoOrd$.

\subsection{Incremental construction}

For large projects, rebuilding $\IG$ from scratch after a single file
edit is expensive.  \texttt{ImportGraphBuilder} supports incremental
updates: given a changed-file set $\Delta \subseteq V$, it re-parses
only the files in $\Delta$ and splices the new edges into the existing
\texttt{PackageFixedPoint}.  Unchanged nodes and edges are reused,
giving $O(|\Delta| \cdot d_{\max})$ incremental cost where $d_{\max}$
is the maximum in-degree.

% ══════════════════════════════════════════════════════════════════════
\section{Dependency Analysis}
\label{sec:analysis}
% ══════════════════════════════════════════════════════════════════════

\subsection{Tarjan's SCC algorithm}

A \emph{strongly connected component} (SCC) of $\IG$ is a maximal subset
$S \subseteq V$ such that every node in $S$ is reachable from every
other node along directed edges.  Cycles manifest as SCCs with
$|S| \geq 2$; singleton SCCs $|S| = 1$ are acyclic nodes.

\DA{} runs Tarjan's depth-first-search-based SCC algorithm
\cite{tarjan}, which discovers all SCCs in $O(|V| + |E|)$ time in a
single pass.

\begin{lstlisting}[style=jugeo-python,
  caption={SCC detection in \texttt{CircularImportDetector}.},
  label={lst:scc}]
def detect_sccs(self, graph: ImportGraph) -> list[SCC]:
    idx, stack, lowlink = [0], [], {}
    index, on_stack, sccs = {}, {}, []

    def strongconnect(v: str) -> None:
        index[v] = lowlink[v] = idx[0]; idx[0] += 1
        stack.append(v); on_stack[v] = True
        for w in graph.successors(v):
            if w not in index:
                strongconnect(w)
                lowlink[v] = min(lowlink[v], lowlink[w])
            elif on_stack.get(w, False):
                lowlink[v] = min(lowlink[v], index[w])
        if lowlink[v] == index[v]:
            comp: set[str] = set()
            while True:
                w = stack.pop(); on_stack[w] = False
                comp.add(w)
                if w == v: break
            sccs.append(SCC(members=frozenset(comp),
                            is_cyclic=len(comp) > 1))

    for v in graph.modules:
        if v not in index:
            strongconnect(v)
    return sccs
\end{lstlisting}

\subsection{Condensation graph}

The \emph{condensation} $\CondGraph$ collapses each SCC to a single
super-node.

\begin{definition}[Condensation]
\label{def:condensation}
Given $\IG = (V, E)$ with SCC partition $\{S_1, \ldots, S_k\}$, the
\emph{condensation} $\CondGraph$ has
\begin{itemize}
  \item $V(\CondGraph) = \{\hat{S}_1, \ldots, \hat{S}_k\}$;
  \item $(\hat{S}_i, \hat{S}_j) \in E(\CondGraph)$ iff
        $\exists\, s \in S_i,\, t \in S_j$ with
        $(s, t) \in E(\IG)$ and $S_i \neq S_j$.
\end{itemize}
\end{definition}

\begin{proposition}[DAG property of $\CondGraph$]
\label{prop:dag}
The condensation $\CondGraph$ is a directed acyclic graph.
\end{proposition}

\begin{proof}
Suppose $\CondGraph$ contains a cycle
$\hat{S}_{i_1} \to \hat{S}_{i_2} \to \cdots \to \hat{S}_{i_1}$.
Then there exist nodes $s_{i_1} \in S_{i_1}, s_{i_2} \in S_{i_2},
\ldots$ forming a cycle in $\IG$ passing through all $S_{i_j}$.
Since nodes in all $S_{i_j}$ are mutually reachable, the union
$\bigcup_j S_{i_j}$ would form a single SCC, contradicting the
assumption that $S_{i_j}$ are distinct SCCs.
\end{proof}

\subsection{Topological ordering via Kahn's algorithm}

Since $\CondGraph$ is a DAG, it admits a topological ordering.

\begin{definition}[Topological order]
\label{def:topo-order}
A \emph{topological order} on DAG $(V, E)$ is a linear order
$<_{\TopoOrd}$ on $V$ such that
\[
  (u, v) \in E \;\Longrightarrow\; u <_{\TopoOrd} v.
\]
For the module graph: if $M$ imports $N$, then $N <_{\TopoOrd} M$,
\ie dependencies come \emph{before} dependents in $\TopoOrd$.
\end{definition}

\DA{} computes $\TopoOrd$ on $\CondGraph$ using Kahn's BFS
\cite{kahn}.  Kahn's algorithm also serves as an acyclicity witness:
if it consumes fewer than $|V(\CondGraph)|$ nodes, a cycle remained—but
since $\CondGraph$ is a DAG (Proposition~\ref{prop:dag}), this cannot
happen.

\begin{lstlisting}[style=jugeo-python,
  caption={Kahn's algorithm for topological order in \DA.},
  label={lst:kahn}]
def topological_order(self, dag: CondensationGraph) -> list[SCC]:
    in_deg = {s: 0 for s in dag.nodes}
    for s in dag.nodes:
        for t in dag.successors(s):
            in_deg[t] += 1
    queue = deque(s for s in dag.nodes if in_deg[s] == 0)
    order: list[SCC] = []
    while queue:
        s = queue.popleft()
        order.append(s)
        for t in dag.successors(s):
            in_deg[t] -= 1
            if in_deg[t] == 0:
                queue.append(t)
    assert len(order) == len(dag.nodes), "BUG: condensation has cycle"
    return order
\end{lstlisting}

% ══════════════════════════════════════════════════════════════════════
\section{Site from Imports}
\label{sec:site}
% ══════════════════════════════════════════════════════════════════════

\subsection{Modules as coordinate objects}

Following the Judgment Geometry framework (Paper~1 \cite{paper01}),
every import graph $\IG$ gives rise to a semantic site $\ModSite$.

\begin{construction}[Module site]
\label{constr:module-site}
Given a well-formed import graph $\IG = (V, E)$, define the
\emph{module site} $\ModSite = (\MdCat, J_{\IG})$ where:
\begin{enumerate}[label=(\alph*)]
  \item \textbf{Category $\MdCat$.}
        Objects are modules $M \in V$.  A morphism $M \to N$ exists
        iff $(M, N) \in E$ (an import edge).  Composition is
        transitivity: $M \to N \to P$ exists whenever $M$ imports $N$
        and $N$ imports $P$.  Identity morphisms come from the
        reflexive closure of reachability.

  \item \textbf{Grothendieck topology $J_{\IG}$.}
        A \emph{covering sieve} on $M$ is generated by the
        \emph{import covering family}:
        \[
          \mathrm{Cov}(M) \coloneqq
          \bigl\{f_N \colon M \to N \;\big|\; N \in \Dep(M)\bigr\}.
        \]
        A sieve $S$ on $M$ covers $M$ iff
        $S \supseteq \mathrm{Cov}(M)$.
\end{enumerate}
\end{construction}

\begin{remark}
In site-theoretic language, each module $M$ carries the judgment
8-tuple
$\J_M = \Jud{M}{\varphi_M}{\carrier_M}{\evid_M}
             {\oblig_M}{\obstr_M}{\trust_M}{\prov_M}$
where $M$ itself is the coordinate $\coord = M$, import edges
$M \to N$ are restriction morphisms
$\res_{M,N} \colon \J_N \to \J_M$ pulling back evidence and
obligations from $N$ to $M$.
\end{remark}

\subsection{The verification sheaf}

\begin{definition}[Verification sheaf]
\label{def:verif-sheaf}
The \emph{verification sheaf} $\mathcal{V}$ assigns to each module $M$
the set of valid judgment tuples at $M$:
\[
  \mathcal{V}(M) = \bigl\{\J_M \mid \trust(\J_M) \tgeq \Tunver\bigr\},
\]
with restriction maps
$\res_{M,N} \colon \mathcal{V}(N) \to \mathcal{V}(M)$
induced by the import edge $(M, N)$.
\end{definition}

\begin{proposition}[Descent along imports]
\label{prop:descent}
The verification sheaf $\mathcal{V}$ satisfies the gluing condition
on $\ModSite$: given a compatible family of local sections
$\{j_N \in \mathcal{V}(N) \mid N \in \Dep(M)\}$, there exists a
unique global section $j_M \in \mathcal{V}(M)$ whose restriction to
each $N$ is $j_N$.
\end{proposition}

\begin{proof}
Compatibility means: for any two imports $N_1, N_2 \in \Dep(M)$ sharing
a common transitive dependency $P$, the sections $j_{N_1}$ and $j_{N_2}$
agree on $P$ after restriction.  Because $M$'s verification is
determined by its imports, the unique glueing is the judgment
$j_M$ built by aggregating the evidence
$\bigoplus_{N \in \Dep(M)} \res_{M,N}(j_N)$.
Uniqueness follows from the fact that the import covering sieve
$\mathrm{Cov}(M)$ is generating.
\end{proof}

\subsection{The \v{C}ech complex of the module site}

The \emph{\v{C}ech complex} $\Cech^{\bullet}(\mathrm{Cov}(M), \mathcal{V})$
computes the cohomological obstructions to gluing.
\begin{itemize}
  \item $\Cech^0$ — sections over each $N \in \Dep(M)$: the local
        verified judgments.
  \item $\Cech^1$ — sections over pairwise overlaps
        $N_1 \cap_M N_2 \coloneqq$ the set of modules in
        $\Dep(N_1) \cap \Dep(N_2)$: compatibility conditions.
  \item $H^1 \neq 0$ signals a gluing obstruction, which arises
        precisely when there is a cycle
        (cf.\ \S\ref{sec:cycles}, Definition~\ref{def:cycle-obstruction}).
\end{itemize}

\subsection{Integration with the 8-tuple judgment}

When module $M$ is verified in the context of $\ModSite$, the
eight components of the JuGeo judgment tuple
$\Jud{c}{\prop}{\carrier}{\evid}{\oblig}{\obstr}{\trust}{\prov}$
are instantiated as follows:

\begin{center}
\begin{tabular}{lll}
\toprule
\textbf{Component} & \textbf{Symbol} & \textbf{Instantiation at module $M$} \\
\midrule
Coordinate  & $\coord$  & Module name $M$ \\
Proposition & $\prop$   & Exported specification of $M$ \\
Carrier     & $\carrier$& Python module object (runtime) \\
Evidence    & $\evid$   & Static analysis + test + proof artefacts \\
Obligations & $\oblig$  & Unresolved import stubs, missing types \\
Obstructions& $\obstr$  & Import cycles, shadowing (\S\ref{sec:cycles}) \\
Trust       & $\trust$  & Level assigned by \DA{} \\
Provenance  & $\prov$   & File path, git commit, analysis timestamp \\
\bottomrule
\end{tabular}
\end{center}

% ══════════════════════════════════════════════════════════════════════
\section{Verification Ordering}
\label{sec:ordering}
% ══════════════════════════════════════════════════════════════════════

\subsection{Verification context and completeness}

\begin{definition}[Verification context]
\label{def:verif-context}
A \emph{verification context} $\VerifCtx$ is a partial function
\[
  \VerifCtx \colon V \rightharpoonup
  \{\mathsf{verified},\, \mathsf{failed},\, \mathsf{unverified}\}
\]
assigning a verification outcome to every processed module.
$\VerifCtx$ is \emph{complete} if $\mathrm{dom}(\VerifCtx) = V$.
\end{definition}

\begin{definition}[Valid global verification]
\label{def:valid-global}
A complete verification context $\VerifCtx$ is a
\emph{valid global verification} of $\IG$ if
\[
  \forall M \in V,\quad \VerifCtx(M) = \mathsf{verified}.
\]
\end{definition}

\subsection{Dependency-ready modules}

At each step of bottom-up traversal, a module $M$ is
\emph{ready for verification} if every direct dependency
$N \in \Dep(M)$ satisfies $\VerifCtx(N) = \mathsf{verified}$.

\begin{lemma}[Readiness in topological order]
\label{lem:readiness}
If modules are processed in order $\TopoOrd$ and every processed
module receives $\mathsf{verified}$, then at the moment $M$ is
processed, every $N \in \Dep(M)$ is already verified.
\end{lemma}

\begin{proof}
Since $(M, N) \in E$ implies $N <_{\TopoOrd} M$, and modules are
processed in the order $<_{\TopoOrd}$, module $N$ is processed before
$M$.  By assumption every processed module is $\mathsf{verified}$,
so $\VerifCtx(N) = \mathsf{verified}$.
\end{proof}

\subsection{Verification schedule}

\DA{} emits the verification schedule as an ordered list of
\emph{verification units}, where each unit is either a singleton
module (acyclic SCC) or an entire SCC bundle (cyclic SCC).

\begin{lstlisting}[style=jugeo-python,
  caption={Generating the verification schedule in \DA.},
  label={lst:schedule}]
def verification_schedule(
        self, graph: ImportGraph) -> list[VerifUnit]:
    sccs  = self.detect_sccs(graph)
    cond  = self.condense(graph, sccs)
    topo  = self.topological_order(cond)
    return [
        VerifUnit(
            scc=s,
            deps=[t for t in topo
                  if cond.has_edge(t.id, s.id)])
        for s in topo
    ]
\end{lstlisting}

The returned list is consumed by the JuGeo verification loop, which
calls \VERIFY{} on each unit in order.  Because \DA{} guarantees
topological ordering, the \VERIFY{} move for unit $u$ may freely
appeal to the verified judgments of all prior units.

% ══════════════════════════════════════════════════════════════════════
\section{Handling Cycles}
\label{sec:cycles}
% ══════════════════════════════════════════════════════════════════════

\subsection{Cyclic imports in Python}

Python permits mutually recursive imports.  A canonical two-module
cycle is:

\begin{lstlisting}[style=jugeo-python,
  caption={Mutual import cycle between \texttt{a.py} and \texttt{b.py}.},
  label={lst:cycle}]
# a.py
from b import B_func   # a imports b

# b.py
from a import A_func   # b imports a
\end{lstlisting}

When Python executes \texttt{import a}, it begins executing
\texttt{a.py}, which triggers \texttt{import b}, which triggers
\texttt{import a} again—but \texttt{a}'s module object is only
partially initialized.  Python returns the partial object, which can
raise \texttt{AttributeError} at runtime.

\subsection{Cycles as site obstructions}

In the site-theoretic picture, a cycle
$M_1 \to M_2 \to \cdots \to M_k \to M_1$
creates a non-trivial class in $H^1(\ModSite, \mathcal{V})$: the local
sections over $\{M_1, \ldots, M_k\}$ cannot be glued without
circularity, because each module's behavior refers to another module
not yet fully defined.

\begin{definition}[Cycle obstruction]
\label{def:cycle-obstruction}
For a cyclic SCC $S = \{M_1, \ldots, M_k\}$, the
\emph{cycle obstruction} is the record
\[
  \CyclObs_S =
  \bigl(\{M_1,\ldots,M_k\},\;
        \mathsf{CycleType},\;
        \obstr_{\mathrm{cycle}}\bigr)
\]
stored in $\obstr(\J_S)$ of the bundled judgment for $S$.
\end{definition}

\subsection{SCC bundling}

\JuGeo{} resolves a cyclic SCC $S$ by bundling: the whole SCC is
treated as a single verification unit with a joint judgment tuple
$\J_S$.

\begin{definition}[Bundle judgment]
\label{def:bundle}
For a cyclic SCC $S = \{M_1, \ldots, M_k\}$, the
\emph{bundle judgment} is
\[
  \J_S = \Jud{S}
             {\bigwedge_{i=1}^k \varphi_{M_i}}
             {\prod_{i=1}^k \carrier_{M_i}}
             {\bigoplus_{i=1}^k \evid_{M_i}}
             {\bigcup_{i=1}^k \oblig_{M_i}}
             {\CyclObs_S \cup \bigcup_{i=1}^k \obstr_{M_i}}
             {\bigwedge_{i=1}^k \trust_{M_i}}
             {\mathrm{bundle}(M_1,\ldots,M_k)},
\]
where $\bigwedge_i \trust_{M_i}$ is the trust-algebra meet
(Paper~4, \S4.2 \cite{paper04}).
\end{definition}

\begin{proposition}[Bundle judgment soundness]
\label{prop:bundle-sound}
If every $M_i \in S$ passes local verification with trust
$\trust_i \tgeq \Tunver$, then
$\trust(\J_S) = \bigwedge_i \trust_i \tgeq \Tunver$.
\end{proposition}

\begin{proof}
The meet of finitely many elements that are all $\tgeq \Tunver$ is
also $\tgeq \Tunver$, by the downward-closure property of the trust
lattice (Paper~4, Theorem~4.3 \cite{paper04}).
\end{proof}

\subsection{SMT encoding of cycle detection}

\DA{} encodes acyclicity of $\IG$ as a \QFLIA{} constraint system.
For each module $M_i \in V$, a rank variable $r_{M_i} \in \NN$ is
introduced.  For each import edge $(M, N) \in E$, the constraint
$r_N < r_M$ is added.  The existence of a cycle corresponds to
unsatisfiability of this system.

The \JuGeo{} SMT dispatcher (Paper~5 \cite{paper05}) sends this query
to Z3 or CVC5.  An \textsf{UNSAT} result is stored in
$\evid(\J_S)$ with $\trust = \Tsolver$.  This gives the obstruction
record a machine-checked certificate: the cycle is not merely
detected heuristically but witnessed by a proof of unsatisfiability.

% ══════════════════════════════════════════════════════════════════════
\section{Theorem: Compositionality}
\label{sec:theorem}
% ══════════════════════════════════════════════════════════════════════

\subsection{Statement and proof}

\begin{theorem}[Compositionality]
\label{thm:compositionality}
Let $\IG = (V, E)$ be a well-formed import graph and let
$\TopoOrd = (M_1, M_2, \ldots, M_n)$ be any topological ordering of
$V$ (\ie $M_i <_{\TopoOrd} M_j$ implies $i < j$).
Let $\LocalVerif \colon V \to \{\mathsf{verified}, \mathsf{failed}\}$
be a local verifier.
Define the verification context
$\VerifCtx_{\TopoOrd}(M_i) \coloneqq \LocalVerif(M_i)$.
Then $\VerifCtx_{\TopoOrd}$ is a valid global verification of $\IG$
if and only if $\LocalVerif(M) = \mathsf{verified}$ for all $M \in V$.
\end{theorem}

\begin{proof}
$(\Rightarrow)$: If $\VerifCtx_{\TopoOrd}$ is a valid global
verification, then by Definition~\ref{def:valid-global},
$\VerifCtx_{\TopoOrd}(M) = \mathsf{verified}$ for all $M \in V$.
Since $\VerifCtx_{\TopoOrd}(M) = \LocalVerif(M)$ by construction,
$\LocalVerif(M) = \mathsf{verified}$ for all $M$.

$(\Leftarrow)$: Assume $\LocalVerif(M) = \mathsf{verified}$ for all
$M \in V$.  Then $\VerifCtx_{\TopoOrd}(M) = \mathsf{verified}$ for
all $M$ by construction.  The context covers $V$ because $\TopoOrd$
is a permutation of $V$.  Hence $\VerifCtx_{\TopoOrd}$ is complete and
all-verified, satisfying Definition~\ref{def:valid-global}.
\end{proof}

\begin{remark}
Theorem~\ref{thm:compositionality} is stated for the acyclic case.
For cyclic imports, the SCC condensation $\CondGraph$ is acyclic by
Proposition~\ref{prop:dag}, and the theorem applies with SCC bundle
judgments (Definition~\ref{def:bundle}) replacing individual modules.
\end{remark}

\subsection{Dependency-awareness corollary}

\begin{corollary}[Dependency-aware soundness]
\label{cor:dep-aware}
Under the hypotheses of Theorem~\ref{thm:compositionality}, if
$\LocalVerif(N) = \mathsf{verified}$ for all $N \in \Dep^*(M)$, then
when $\LocalVerif(M)$ is invoked, the restriction of
$\VerifCtx_{\TopoOrd}$ to $\Dep^*(M)$ is complete and all-verified.
\end{corollary}

\begin{proof}
By Lemma~\ref{lem:readiness} and the hypothesis on $\LocalVerif$.
\end{proof}

\subsection{Lean formalization}

The proof is mechanized in \texttt{Paper26\_Import\-Graph.lean} inside
namespace \texttt{JudgmentGeometry.ImportGraph}.  The key definitions
and the main theorem are sketched below.

\begin{verbatim}
structure ModuleId  where name : String
  deriving DecidableEq, BEq, Repr

structure ImportEdge where src dst : ModuleId
  deriving DecidableEq, Repr

structure ImportGraph where
  modules : List ModuleId
  edges   : List ImportEdge

-- Position in list (length if absent, treated as infinity)
def posInList (m : ModuleId) (l : List ModuleId) : Nat :=
  l.indexOf m

-- A list `order` is a valid topological order for `g`
def ImportGraph.isTopoOrder (g : ImportGraph)
    (order : List ModuleId) : Prop :=
  forall e : ImportEdge, e in g.edges ->
    posInList e.dst order < posInList e.src order

-- Build a verification context by processing modules in order
def buildVerifContext (order : List ModuleId)
    (verify : ModuleId -> VerifStatus) : VerifContext :=
  order.map (fun m => { module := m, status := verify m })

-- Compositionality (Theorem 7.1)
theorem compositionality
    (g      : ImportGraph)
    (order  : List ModuleId)
    (horder : g.isTopoOrder order)
    (hcovs  : forall m in g.modules, m in order)
    (verify : ModuleId -> VerifStatus)
    (hall   : forall m in order, verify m = VerifStatus.verified) :
    validGlobalVerif (buildVerifContext order verify) g
\end{verbatim}

The proof splits into two helper lemmas:
\begin{enumerate}[label=(\roman*)]
  \item \texttt{buildVerifContext\_allVerified}: the context is
        all-verified whenever \texttt{hall} holds;
  \item \texttt{buildVerifContext\_covers}: the context covers every
        module in \texttt{order}.
\end{enumerate}
Both are proved without \texttt{sorry}.  The full file also
formalizes SCC properties (\texttt{SCC.isTrivial},
\texttt{SCC.isCyclic} and their exclusion theorems),
the converse direction
(\texttt{globalVerif\_implies\_local}), and the negative result
(\texttt{empty\_ctx\_not\_valid}).

% ══════════════════════════════════════════════════════════════════════
\section{Experiments}
\label{sec:experiments}
% ══════════════════════════════════════════════════════════════════════

\subsection{Benchmark setup}

We evaluate \DA{} on the 10 JuGeo benchmark programs summarised in
\texttt{data-paper26.tex}.  The audited measurements therefore concern
the repository's benchmark corpus only; we do not claim validated
results for external projects such as NumPy or Django in this version.

\subsection{Graph construction and analysis time}

\begin{table}[H]
\centering
\caption{Import graph construction metrics (JuGeo benchmark programs).}
\label{tab:times}
\begin{tabular}{lr}
\toprule
\textbf{Metric} & \textbf{Value} \\
\midrule
Programs analysed       & \ppTwentysixTotalPrograms \\
Mean coordinates/site   & \ppTwentysixMeanCoords \\
Mean morphisms/site     & \ppTwentysixMeanMorphisms \\
Mean site build time    & \ppTwentysixMeanBuildMs\,ms \\
Max site build time     & \ppTwentysixMaxBuildMs\,ms \\
Total functions indexed & \ppTwentysixTotalFunctions \\
Total classes indexed   & \ppTwentysixTotalClasses \\
\bottomrule
\end{tabular}
\end{table}

Construction time is dominated by AST traversal; the site build
overhead of \ppTwentysixMeanBuildMs\,ms per program is negligible
relative to the end-to-end verification time of \ppTwentysixUnordMeanTime{}.

\subsection{Cycle detection}

\begin{table}[H]
\centering
\caption{Cycle detection and obstruction results.}
\label{tab:cycles}
\begin{tabular}{lr}
\toprule
\textbf{Metric} & \textbf{Value} \\
\midrule
Programs verified          & \ppTwentysixVerifiedCount \\
Total obstructions found   & \ppTwentysixTotalObstructions \\
Overall verification rate  & \ppTwentysixAccuracy \\
Trust obligations discharged & \ppTwentysixTrustDischarged \\
\bottomrule
\end{tabular}
\end{table}

The JuGeo benchmark suite is obstruction-free by design: all
\ppTwentysixTotalPrograms{} programs verify with zero cyclic import
obstructions.  \DA{} confirms this by verifying rank constraints
with trust level $\trust = \Tsolver$.

\subsection{Verification ordering speedup}

We compare two strategies on JuGeo core:

\begin{enumerate}[label=(\roman*)]
  \item \textbf{Unordered}: modules verified in arbitrary
        file-system traversal order.
  \item \textbf{Topological}: modules verified in the order produced
        by \DA{}.
\end{enumerate}

\begin{table}[H]
\centering
\caption{Verification time comparison.}
\label{tab:speedup}
\begin{tabular}{lrr}
\toprule
\textbf{Strategy} & \textbf{Re-verifications} & \textbf{Time (ms)} \\
\midrule
Unordered    & \ppTwentysixUnordMeanReverif & \ppTwentysixUnordMeanTime \\
Topological  & \ppTwentysixTopoReverif & \ppTwentysixTopoMeanTime \\
\bottomrule
\end{tabular}
\end{table}

Unordered verification triggers re-verifications because dependents
are processed before dependencies, requiring back-propagation.
Topological verification eliminates those re-verifications entirely, but
in the current benchmark it is slightly slower in wall-clock time
(4893.54\,ms vs.\ 4725.16\,ms mean).  The measured benefit is therefore
reduced redundant work rather than a net speedup on this dataset.

\subsection{Site construction integration}

The module site $\ModSite$ produced by \DA{} is passed to
\texttt{jugeo.geometry.site.build\_site}, which constructs the full
Grothendieck site object consumed by the semantic move layer
(Paper~6 \cite{paper06}).  The integration bridge
\texttt{ImportsPackageFixedPointsBridge} converts SCCs, obstructions,
and trust annotations into JuGeo judgment tuples, enabling downstream
use of \VERIFY{}, \GLUE{}, and \REPAIR{} moves.

% ══════════════════════════════════════════════════════════════════════
\section{Conclusion}
\label{sec:conclusion}
% ══════════════════════════════════════════════════════════════════════

We have shown that the import graph of a Python project is not merely
organizational metadata but a mathematical structure—a Grothendieck
site—that directly shapes the verification strategy.  \DA{} extracts
this structure from Python AST, detects cycles via Tarjan's SCC
algorithm, and emits a topological verification order via Kahn's
algorithm.

The import graph remains useful as a dependency model for bottom-up
verification ordering.
When a module is processed, all its dependencies already carry
verified judgments, so local verification may freely cite them.
In the current benchmark, cyclic-import handling does not surface any
obstructions: all 10 benchmark programs are obstruction-free.

\paragraph{Future directions.}
\begin{enumerate}[label=(\roman*)]
  \item \textbf{Incremental re-verification.}
        After a single file edit, only the affected SCC and its
        transitive dependents need re-verification.
  \item \textbf{Parallel verification.}
        Independent SCCs in $\CondGraph$ can be verified
        concurrently without violating topological ordering.
  \item \textbf{Cross-language extension.}
        Paper~30's motivic transfer \cite{paper30} can propagate
        verified judgments across language boundaries when a Python
        module wraps a compiled extension.
\end{enumerate}

% ══════════════════════════════════════════════════════════════════════
\begin{thebibliography}{99}

\bibitem{paper01}
H.~Young,
\emph{Semantic Sites for Program Verification},
Judgment Geometry Series, Paper~1, 2025.

\bibitem{paper02}
H.~Young,
\emph{Judgment Algebra},
Judgment Geometry Series, Paper~2, 2025.

\bibitem{paper03}
H.~Young,
\emph{Descent Obstructions in Verification},
Judgment Geometry Series, Paper~3, 2025.

\bibitem{paper04}
H.~Young,
\emph{Trust Algebra for Multi-Tier Verification},
Judgment Geometry Series, Paper~4, 2025.

\bibitem{paper05}
H.~Young,
\emph{SMT Fragment Dispatch},
Judgment Geometry Series, Paper~5, 2025.

\bibitem{paper06}
H.~Young,
\emph{Semantic Moves},
Judgment Geometry Series, Paper~6, 2025.

\bibitem{paper07}
H.~Young,
\emph{Python Effects in the Judgment Framework},
Judgment Geometry Series, Paper~7, 2025.

\bibitem{paper30}
H.~Young,
\emph{Motivic Measures for Cross-Language Verification},
Judgment Geometry Series, Paper~30, 2025.

\bibitem{tarjan}
R.~E.~Tarjan,
\emph{Depth-first search and linear graph algorithms},
SIAM Journal on Computing, 1(2):146--160, 1972.

\bibitem{kahn}
A.~B.~Kahn,
\emph{Topological sorting of large networks},
Communications of the ACM, 5(11):558--562, 1962.

\bibitem{dafny}
K.~R.~M.~Leino,
\emph{Dafny: An automatic program verifier for functional correctness},
LPAR 2010: Logic for Programming, Artificial Intelligence,
and Reasoning, Springer, 2010.

\bibitem{sga4}
M.~Artin, A.~Grothendieck, J.-L.~Verdier,
\emph{Th\'eorie des Topos et Cohomologie \'Etale
      des Sch\'emas (SGA~4)},
Lecture Notes in Mathematics 269--270, Springer-Verlag, 1972.

\end{thebibliography}

\end{document}
